\documentclass[12pt,letterpaper]{article}
\usepackage[utf8]{inputenc}
\usepackage[spanish]{babel}
\usepackage[margin=2.5cm]{geometry}
\usepackage{amsmath}
\usepackage{amsfonts}
\usepackage{amssymb}
\usepackage{array}

\begin{document}

\begin{center}
    \textbf{\Large Formulario de la Unidad II}\\
    \vspace{0.2cm}
    \textbf{Física para Ciencias de la Salud (005-1713)}\\
    \vspace{0.2cm}
    \textit{Cinemática, Dinámica, Trabajo, Energía y Potencia}
\end{center}

\vspace{0.5cm}

\begin{table}[h!]
    \centering
    \renewcommand{\arraystretch}{2} % Aumenta el espaciado vertical de las filas para mayor legibilidad
    \begin{tabular}{|>{\raggedright\arraybackslash}p{4cm}|>{\centering\arraybackslash}p{5.5cm}|>{\raggedright\arraybackslash}p{5cm}|}
        \hline
        \textbf{Concepto} & \textbf{Fórmula} & \textbf{Variables} \\
        \hline
        \textbf{Movimiento Rectilíneo Uniforme (MRU)} & $v = \dfrac{d}{t}$ & $v$: velocidad \newline $d$: distancia \newline $t$: tiempo \\
        \hline
        \textbf{Aceleración Media} & $a = \dfrac{v_f - v_i}{t}$ & $a$: aceleración \newline $v_f, v_i$: vel. final e inicial \\
        \hline
        \textbf{Cinemática (MRUA)} \newline (Con aceleración constante) & $v_f = v_i + a \cdot t$ \newline $d = v_i \cdot t + \frac{1}{2}a \cdot t^2$ \newline $v_f^2 = v_i^2 + 2a \cdot d$ & $d$: desplazamiento \newline $a$: aceleración \\
        \hline
        \textbf{Segunda Ley de Newton} & $\sum \vec{F} = m \cdot \vec{a}$ & $F$: fuerza neta (N) \newline $m$: masa (kg) \\
        \hline
        \textbf{Peso} \newline (Fuerza gravitatoria) & $W = m \cdot g$ & $W$: peso (N) \newline $g$: gravedad ($\approx 9.8 \, \text{m/s}^2$) \\
        \hline
        \textbf{Fuerza Centrípeta} & $F_c = m \cdot a_c = m \cdot \dfrac{v^2}{r}$ & $F_c$: fuerza centrípeta \newline $a_c$: ac. centrípeta \newline $r$: radio de giro \\
        \hline
        \textbf{Trabajo Mecánico} & $W = F \cdot d \cdot \cos(\theta)$ & $W$: trabajo (Joules, J) \newline $\theta$: ángulo entre $F$ y $d$ \\
        \hline
        \textbf{Energía Cinética} & $E_c = \frac{1}{2} m \cdot v^2$ & $E_c$: energía cinética (J) \newline $v$: velocidad \\
        \hline
        \textbf{Energía Potencial} \newline (Gravitatoria) & $E_p = m \cdot g \cdot h$ & $E_p$: energía potencial (J) \newline $h$: altura (m) \\
        \hline
        \textbf{Potencia Mecánica} & $P = \dfrac{W}{t} \quad \text{o} \quad P = F \cdot v$ & $P$: potencia (Watts, W) \newline $t$: tiempo (s) \\
        \hline
    \end{tabular}
    \vspace{0.3cm}
    \caption{Fórmulas fundamentales para la resolución de problemas (Unidad II).}
    \label{tab:formulas_u2}
\end{table}

\vspace{0.5cm}
\noindent \textbf{Nota para conversiones comunes:}
\begin{itemize}
    \item Fuerza: $1 \text{ N} = 1 \text{ kg} \cdot \text{m/s}^2$
    \item Trabajo/Energía: $1 \text{ Joule (J)} = 1 \text{ N} \cdot \text{m}$
    \item Potencia: $1 \text{ Watt (W)} = 1 \text{ J/s}$
\end{itemize}

\end{document}